\documentclass[uplatex,a4paper,12pt]{jsarticle}
\usepackage{amsmath,amssymb}
\usepackage{enumitem}

\begin{document}

\title{ニューラルネットワーク基礎 確認テスト}
\author{}
\date{}
\maketitle

\section*{注意}
\begin{itemize}
\item 試験時間は5分とする．
\item 全10問，各10点，合計100点とする．
\item 60点以上を合格とする．
\end{itemize}

\section*{問題}

\begin{enumerate}

\item 次のうち，多層パーセプトロン（MLP）の説明として最も適切なものを1つ選べ．
\begin{enumerate}[label=(\alph*)]
\item 線形回帰のみを行うモデル
\item パーセプトロンを多層化したニューラルネットワーク
\item データを暗号化するモデル
\end{enumerate}

\item 次のうち，線形分離不可能な問題として知られているものを1つ選べ．
\begin{enumerate}[label=(\alph*)]
\item AND問題
\item OR問題
\item XOR問題
\end{enumerate}

\item 次の式
\[
y=g(\mathbf{w}^\top \mathbf{x}+b)
\]
において，$g$ に相当するものを1つ選べ．
\begin{enumerate}[label=(\alph*)]
\item 誤差関数
\item 活性化関数
\item 学習率
\end{enumerate}

\item 次のうち，ReLU関数として正しいものを1つ選べ．
\begin{enumerate}[label=(\alph*)]
\item $\mathrm{ReLU}(x)=\max(x,0)$
\item $\mathrm{ReLU}(x)=\dfrac{1}{1+e^{-x}}$
\item $\mathrm{ReLU}(x)=x^2$
\end{enumerate}

\item ニューラルネットワークにおいて，入力層と出力層の間に存在する層を何というか．

\item 次のうち，多値分類問題で出力層によく用いられる活性化関数を1つ選べ．
\begin{enumerate}[label=(\alph*)]
\item ReLU
\item ソフトマックス関数
\item 恒等関数
\end{enumerate}

\item 勾配降下法における学習率 $\eta$ の役割として正しいものを1つ選べ．
\begin{enumerate}[label=(\alph*)]
\item パラメータ更新量を調整する
\item クラス数を決定する
\item データ数を増加させる
\end{enumerate}

\item 次の説明に当てはまる語句を答えよ．

「出力層から順に誤差を伝播させることで，各パラメータの勾配を効率的に計算する手法」

\item 次のうち，ハイパーパラメータに該当するものを1つ選べ．
\begin{enumerate}[label=(\alph*)]
\item 重み $W$
\item バイアス $b$
\item 学習率
\end{enumerate}

\item 次の説明として正しいものを1つ選べ．
\begin{enumerate}[label=(\alph*)]
\item MLPは非線形関数を表現できない
\item 隠れ層と活性化関数によって非線形な関数を近似できる
\item 活性化関数を用いなくても非線形問題を解ける
\end{enumerate}

\end{enumerate}

\end{document}